\begin{enumerate}[label=\textbf{\ansnum.\arabic*}]
  \item $v_1,v_2,w_{11},w_{21}$ have negative gradients, so they increase. $w_{12},w_{22}$ have positive gradients, so they decrease. The prediction $o\approx0.55$ is too small. Hidden neuron $h_1$ is connected to the output with a positive weight, so making $h_1$ larger (larger $w_{11},w_{21}$) raises $o$. Hidden neuron $h_2$ is connected with a negative weight $v_2$, so making $h_2$ \emph{smaller} (smaller $w_{12},w_{22}$) raises $o$. Increasing $v_1$ and $v_2$ also raises $o$ because $h_1,h_2>0$.
  \item $\delta_o=(o-t)\,o(1-o)=0.5452\cdot0.2480\approx0.1352$ and $\pder{L}{v_1}=0.1352\cdot0.6341\approx0.0857$. The gradient is positive, so $v_1$ decreases: the output is too large for the target $0$ and has to be reduced.
  \item All four gradients are $0$, because each is $\delta_j\cdot x_i$ with $x_i=0$: a weight attached to a zero input has no influence on the loss. $v_1$ still receives a gradient, since $h_1=h_2=\sig(0)=0.5$ is not zero.
  \item Each sigmoid layer multiplies the signal by a factor $\sig'(z)\le\tfrac14$ (and by weights). After $n$ layers the factor can be as small as $(1/4)^n$, so the gradients of the early layers become tiny and these layers learn very slowly. This is one reason why ReLU is popular.
  \item With $\eta=0.5$ the loss falls from $0.1034$ to about $0.003$ after $200$ steps and the output approaches $1$. With $\eta=0.01$ the loss only falls to about $0.087$, so learning is very slow. With $\eta=10$ this tiny network still learns (the sigmoid keeps the steps bounded), but in general very large learning rates make the loss jump around or diverge, so $\eta$ has to be chosen with care.
  \item The central difference $\dfrac{L(v_1+h)-L(v_1-h)}{2h}$ with $h=10^{-5}$ gives about $-0.0715$.
  \item $\pder{L}{o}=-0.4548$, $\pder{o}{z_o}=0.2480$, $\pder{z_o}{h_1}=v_1=0.7$, $\pder{h_1}{z_1}=0.2320$, $\pder{z_1}{w_{11}}=x_1=1$. The product is $-0.4548\cdot0.2480\cdot0.7\cdot0.2320\cdot1\approx-0.0183$.
\end{enumerate}
